% !PRINTED 727
\weberpage{727}

%<<CONT
der Formel (9) alle Werthe von $\beta$ auf das Intervall von $0$ bis
$\nu' - 1$ reduciren, so bekommen wir alle Zahlen dieses Intervalles.

Demnach lässt sich die Formel (3) so darstellen:
\weq{10}{
X_1 = -\frac{2}{m}\,(1,\,\Theta^{-1},\,r)\sum_{0,\,\nu'-1}^{\beta}\Theta^\beta\lambda_\beta.
}
Dies haben wir in dem im \S.\,197, (8) angegebenen Ausdrucke
\[
D B = 2^{2^{\varkappa-4}(\varkappa-2)}\prod X_1 = m^{\nu''}2^{-\nu'}\prod X_1
\]
einzusetzen, und dann ergiebt sich, wenn wir
\weq{11}{
U_\Theta = \sum_{0,\,\nu'-1}^{\beta}\Theta^\beta\lambda_\beta
}
setzen, mit Rücksicht auf \S.\,197, (14):
\weq{12}{
D B = \prod^{\Theta} U_\Theta.
}
Dies Product lässt sich nun auf folgende Weise durch eine
Determinante darstellen, die nur die conjugirten Logarithmen
der Einheiten $\tau$ enthält.

Lassen wir den Index $\Theta$ bei $U$ weg, so können wir das
Gleichungssystem ansetzen:
\[
\begin{array}{l}
U \;=\; \lambda_0 \;+\; \lambda_1\Theta \;+\; \lambda_2\Theta^2 \;+\; \cdots \;+\; \lambda_{\nu'-1}\Theta^{\nu'-1} \\[4pt]
\Theta U \;=\; -\lambda_{\nu'-1} \;+\; \lambda_0\Theta \;+\; \lambda_1\Theta^2 \;+\; \cdots \;+\; \lambda_{\nu'-2}\Theta^{\nu'-1} \\[4pt]
\text{\normalfont .\;.\;.\;.\;.\;.\;.\;.\;.\;.\;.\;.\;.\;.\;.\;.\;.\;.\;.\;.\;.\;.\;.\;.\;.\;.\;.\;.\;.\;.} \\[4pt]
\Theta^{\nu'-1} U \;=\; -\lambda_1 \;-\; \lambda_2\Theta \;-\; \lambda_3\Theta^2 \;+\; \cdots \;+\; \lambda_0\Theta^{\nu'-1},
\end{array}
\]
woraus durch Elimination der Potenzen von $\Theta$ folgt:
\weq{13}{
\begin{vmatrix}
\lambda_0 - U, & \lambda_1, & \lambda_2, & \ldots, & \lambda_{\nu'-1} \\
-\lambda_{\nu'-1}, & \lambda_0 - U, & \lambda_1, & \ldots, & \lambda_{\nu'-2} \\
\multicolumn{5}{c}{\text{\normalfont .\;.\;.\;.\;.\;.\;.\;.\;.\;.\;.\;.\;.\;.\;.\;.}} \\
-\lambda_1, & -\lambda_2, & -\lambda_3, & \ldots, & \lambda_0 - U
\end{vmatrix} = 0.
}
% [source note: the first row prints a semicolon-shaped mark after \lambda_0-U; the mathematically required list punctuation is normalized to a comma.]
Dies ist für $U$ eine Gleichung $\nu'^{\text{ten}}$ Grades, deren Wurzeln
die Grössen $U_\Theta$ sind. Ordnet man nach Potenzen von $U$, so
erhält, da $\nu'$ gerade ist, die höchste Potenz $U^{\nu'}$ von $U$ den
Co\"efficienten $1$, und man findet also das in (12) vorkommende
Product der Wurzeln, wenn man $U = 0$ setzt, aus der Determinante (13).

Die so gewonnene Determinante wollen wir etwas anders anordnen,
indem wir die erste Zeile an ihrer Stelle lassen, die
%CONT>>
